%% 校招/实习一页骨架。占位符不是真人；只填候选人确认的事实。
%% xelatex / ctexart / xeCJK，复用上一级 cnresume.sty。
%% 从本目录编译：TEXINPUTS=.:..//: xelatex -halt-on-error resume.tex
%% 复制到 cv/ 后沿用 templates/cn/README.md 的编译命令。
\documentclass[a4paper,fontset=none,zihao=5]{ctexart}
\usepackage{cnresume}
% 照片默认关闭。仅经本人要求时启用：
% \cnphoto{photo.jpg}
\hypersetup{pdftitle={校招实习简历},pdfauthor={姓名},pdfsubject={Resume}}
\begin{document}
\cnheader{[姓名]}{[校招或实习意向] · [城市] · [可到岗时间]}{%
  [电话]\quad\href{mailto:[邮箱]}{[邮箱]}%
}
\cnsection{教育背景}
\cnedu{[入学]-[预计毕业]}{[专业、在读或已获学位]}{[学校]}{[城市]}
\cnprofile{[毕业届别；可选：真实 GPA/满分、排名/总人数、相关课程。未提供则删除。]}

\cnsection{实习经历}
% 无实习时删除本节；不能把课程项目改成任职经历。
\cnrole{[开始]-[结束]}{[实习岗位]}{[实习单位]}{[城市]}
\begin{itemize}
  \item {[具体任务、个人贡献和可核对结果；没有指标不编数字。]}
  \item {[使用的工具、验证方式与复盘。]}
\end{itemize}

\cnsection{项目经历}
\cnproject{[项目名称]}{[开始]-[结束]}{[课程、个人或团队项目；个人角色]}
\begin{itemize}
  \item {[问题、方案、自己完成的部分。]}
  \item {[验证结果、局限和下一步；注明测试环境。]}
\end{itemize}
\cnproject{[项目名称]}{[开始]-[结束]}{[背景与个人角色]}
\begin{itemize}
  \item {[一项可解释的贡献。]}
\end{itemize}

\cnsection{在校经历 / 竞赛 / 奖学金}
\cnproject{[在校经历或竞赛]}{[时间]}{[组织、角色或主办方；可选，无则删除]}
\begin{itemize}
  \item {[个人分工、协作过程、奖项级别；团队奖注明团队及本人贡献。]}
\end{itemize}
\cnprofile{[奖学金名称、授予单位与时间；可选，无则删除。]}

\cnsection{专业技能}
\cnskill{技术：}{[能在面试中解释和演示的技能]}
\cnskill{语言：}{[语言与真实水平]}
\cnskill{实习安排：}{[每周到岗天数、持续月数；仅实习填写]}
\end{document}
